% hori.tex - 横排 clreq 管道（H1）回归测试
% 覆盖：CJK 断点注入、行首行尾禁则（penalty）、中西间距 1/4em 及其例外、
%       符号分离禁则（数字串/数字+单位/货币）、两字宽标点不可拆
\documentclass[12pt]{article}
\usepackage[paperwidth=10cm, paperheight=16cm, margin=1.2cm]{geometry}
% 大陆偏靠式标点（style=mainland 默认）需要 SC 字面：思源宋体标点靠始端、
% 空白在末端；TW-Kai 的标点为台式居中字面，不能用于大陆式测试。
% Script=CJK + Language 必须显式指定：否则 luaotfload 按 latn/dflt 起排，
% 该语言系统的 locl 会把省略号替换为贴基线的西文形（cid63070）；
% hani/ZHS 下省略号保持居中默认形，破折号经 locl 升级为满格中文形。
\usepackage{fontspec}
\setmainfont{Source Han Serif SC}[Script=CJK, Language={Chinese Simplified}]
\usepackage{luatex-cn-hori}
\pagestyle{empty}
\begin{document}

% --- 基本断行与行首行尾禁则：长句在窄版面上多次换行，
%     点号不应出现在行首、开引号不应出现在行尾
天地玄黄，宇宙洪荒。日月盈昃，辰宿列张。寒来暑往，秋收冬藏。闰余成岁，律吕调阳。云腾致雨，露结为霜。金生丽水，玉出昆冈。剑号巨阙，珠称夜光。果珍李柰，菜重芥姜。

子曰：「学而时习之，不亦说乎？有朋自远方来，不亦乐乎？人不知而不愠，不亦君子乎？」又曰：「温故而知新，可以为师矣。」

% --- 中西混排：汉字与西文、数字之间应有约 1/4 字宽的间距
本项目基于 LuaTeX 引擎开发，版本号为 1.18.0，托管在 GitHub 平台上。使用 expl3 编程层与 Lua 脚本协同工作，支持 OpenType 字体特性。

% --- 中西间距例外：点号与夹注号旁不加间距
他说：Hello 世界！括号（English）内外、标点，abc 之后。

% --- 符号分离禁则：数字串、数字+单位、货币不拆行
测量结果为 1234 毫米，占比 95\%，温度 37℃，误差 ±3，价格 ¥1280，共计 100 件样品，编号从 001 到 999，全部合格。

% --- 同上但源码不加空格：中西间距 1/4em 应由管道自动插入
复测结果为5678毫米，占比87\%，温度36℃，误差±2，价格¥999，共计50件样品，标号从010到099，使用LuaTeX编译，全部达标。

% --- 两字宽标点：破折号与省略号
生活就像一盒巧克力——你永远不知道下一颗是什么味道……这就是命运的安排。

\end{document}
